% Part: first-order-logic
% Chapter: axiomatic-deduction
% Section: identity

\documentclass[../../../include/open-logic-section]{subfiles}

\begin{document}

\olfileid{fol}{axd}{ide}

\olsection{\usetoken{P}{derivation} با \usetoken{S}{identity}}

برای گنجاندن $\eq$ در !!{derivation}s، کافی است شِماهای اصل موضوع تازه‌ای
بیفزاییم. تعریف !!{derivation} و $\Proves$ همان می‌ماند؛ فقط اصول موضوع
تازه را نیز مجاز می‌دانیم.

\begin{defn}[اصول موضوع برای !!{identity}]
\begin{align}
\ollabel{ax:id1} & \eq[t][t], \\
\ollabel{ax:id2} & \eq[t_1][t_2] \lif (!B(t_1) \lif
  !B(t_2)),
\end{align}
که در آن $t$، $t_1$ و~$t_2$ همگی ترم بسته‌اند.
\end{defn}

\begin{prop}\ollabel{prop:sound}
  اصول موضوع \olref{ax:id1} و \olref{ax:id2} معتبرند.
\end{prop}

\begin{proof}
  تمرین.
\end{proof}

\begin{prob}
\olref[fol][axd][ide]{prop:sound} را اثبات کنید.
\end{prob}

\begin{prop}\ollabel{prop:iden1}
 $\Gamma \Proves \eq[t][t]$، برای هر ترم بستهٔ $t$ و مجموعهٔ~$\Gamma$.
\end{prop}

\begin{prop}\ollabel{prop:iden2}
  اگر $\Gamma \Proves !A(t_1)$ و $\Gamma \Proves \eq[t_1][t_2]$، آنگاه
  $\Gamma \Proves !A(t_2)$، به شرط بسته‌بودن هر دو ترم~$t_1$ و~$t_2$.
\end{prop}

\begin{proof}
!!{formula} زیر
\[
(\eq[t_1][t_2] \lif (!A(t_1) \lif !A(t_2)))
\]
نمونه‌ای از~\olref{ax:id2} است. نتیجه با دو بار کاربرد~\MP به دست می‌آید.
\end{proof}

\end{document}
